\section{EVALUATION INSTRUMENTS}

\subsection{SUS Questionnaire}

The following ten-item System Usability Scale questionnaire~\cite{brooke1996}
was administered to each participant immediately after completing the task
on each system. Responses were collected on a 5-point Likert scale ranging
from ``Strongly Disagree'' (1) to ``Strongly Agree'' (5).

\begin{longtable}{lp{0.75\linewidth}}
  \caption{System Usability Scale items as administered} \label{tab:sus_items} \\
  \toprule
  \textbf{Item} & \textbf{Statement} \\
  \midrule
  \endfirsthead
  \multicolumn{2}{c}{\tablename\ \thetable{} (continued)} \\
  \toprule
  \textbf{Item} & \textbf{Statement} \\
  \midrule
  \endhead
  \bottomrule
  \endlastfoot
  SUS-1  & I think that I would like to use this system frequently. \\
  SUS-2  & I found the system unnecessarily complex. \\
  SUS-3  & I thought the system was easy to use. \\
  SUS-4  & I think that I would need the support of a technical person to be able to use this system. \\
  SUS-5  & I found the various functions in this system were well integrated. \\
  SUS-6  & I thought there was too much inconsistency in this system. \\
  SUS-7  & I would imagine that most people would learn to use this system very quickly. \\
  SUS-8  & I found the system very cumbersome to use. \\
  SUS-9  & I felt very confident using the system. \\
  SUS-10 & I needed to learn a lot of things before I could get going with this system. \\
\end{longtable}

\noindent
SUS scoring: Odd-numbered items contribute (response $-$ 1); even-numbered
items contribute (5 $-$ response). The sum of all contributions is
multiplied by 2.5 to produce a score on the 0--100 scale.

\subsection{Task Scenario Script}

The following script was read to participants before each testing session.
The same script was used for both the new platform and the WordPress
baseline conditions, with the addition of point (4) for the new platform
only.

\begin{quote}
\itshape
You are a lecturer at VNU-HCM preparing a short interactive video on a
scientific topic for your online course. You have been given a 5-minute
lecture video. Your task is to:

\begin{enumerate}
  \item Upload the provided video file to the system (or paste the
  YouTube URL that has been provided to you).
  \item Add at least three H5P questions at appropriate points in the
  video using different question types.
  \item Export the final H5P package file.
\end{enumerate}

\textit{[New platform only:]} (4) Before adding your own questions, use
the AI analysis feature to generate suggestions from the video transcript
and review them. Accept any suggestions you find suitable.

Please think aloud as you work---tell us what you are trying to do, what
you expect to happen, and what you observe. There are no wrong answers; we
are testing the system, not you. You may stop at any time.
\end{quote}

\subsection{Individual SUS Scores}

Table~\ref{tab:sus_individual} records individual SUS scores for all ten
participants across both conditions.

\begin{table}[H]
  \centering
  \caption{Individual SUS scores by participant}
  \label{tab:sus_individual}
  \begin{tabularx}{\linewidth}{lll}
    \toprule
    \textbf{Participant} & \textbf{New Platform} & \textbf{WordPress Baseline} \\
    \midrule
    P01 & 87.5 & 42.5 \\
    P02 & 82.5 & 35.0 \\
    P03 & 95.0 & 27.5 \\
    P04 & 75.0 & 50.0 \\
    P05 & 90.0 & 20.0 \\
    P06 & 80.0 & 45.0 \\
    P07 & 72.5 & 32.5 \\
    P08 & 85.0 & 55.0 \\
    P09 & 77.5 & 30.0 \\
    P10 & 80.0 & 42.5 \\
    \midrule
    \textbf{Mean} & \textbf{82.5} & \textbf{38.0} \\
    \textbf{SD}   & \textbf{6.8}  & \textbf{12.4} \\
    \bottomrule
  \end{tabularx}
\end{table}
